\documentclass{beamer}

% Theme
\usetheme{Madrid}
\usecolortheme{beaver}

% Metadata
\title{Optimizing YOLOv11 for Edge Deployment}
\subtitle{Quantization and implementation on Raspberry Pi}
\author{Final Year Project Team}
\date{\today}

\begin{document}

% Title Slideå
\begin{frame}
    \titlepage
\end{frame}

% Introduction
\begin{frame}{Introduction}
    \begin{itemize}
        \item \textbf{Objective}: Deploy a License Plate Recognition (LPR) system on a Raspberry Pi.
        \item \textbf{Constraint}: The Raspberry Pi has limited compute (CPU-only inference) and memory compared to desktop GPUs.
        \item \textbf{Solution}: Use a highly efficient object detection model (YOLOv11) and optimize it using \textbf{Quantization}.
    \end{itemize}
\end{frame}

% The Model
\begin{frame}{The Model: YOLOv11}
    \begin{itemize}
        \item \textbf{Architecture}: Ultralytics YOLOv11 (segmentation/detection).
        \item \textbf{Original Format}: PyTorch (.pt), FP32 (Floating Point 32-bit).
        \item \textbf{Issue}: FP32 models are heavy (~20MB+) and slow on ARM CPUs.
        \item \textbf{Optimization Target}: INT8 (Integer 8-bit), reducing size by $\approx 4\times$ and effectively utilizing ARM NEON instructions.
    \end{itemize}
\end{frame}

% Quantization Workflow: Preparation
\begin{frame}{Quantization Workflow: Preparation}
    \begin{enumerate}
        \item \textbf{Environment Setup}: Installed `onnx`, `onnxruntime`, and `onnxruntime-tools`.
        \item \textbf{Model Validation}: Loaded the original FP32 model (`license-plate-finetune-v1l.onnx`) and verified its structural integrity.
        \item \textbf{Input Identification}: Identified the serving input node `'images'` accepting tensor shape $(1, 3, 640, 640)$.
    \end{enumerate}
\end{frame}

% Quantization Workflow: Calibration
\begin{frame}{Quantization Workflow: Calibration}
    \begin{itemize}
        \item \textbf{Why Calibration?} Static quantization requires running representative data through the model to determine dynamic ranges (min/max) for activations.
        \item \textbf{Calibration Reader}: Implemented a `LicensePlateCalibrationReader` class to:
        \begin{itemize}
            \item Preprocess images (Resize, Normalize).
            \item Reorder channels (HWC $\to$ CHW).
            \item Batch data for the quantizer.
        \end{itemize}
    \end{itemize}
\end{frame}

% Quantization Process
\begin{frame}{Process: Static INT8 Quantization}
    \begin{itemize}
        \item We used the ONNX Runtime Quantization API.
        \item \textbf{Method}: Static Quantization.
        \item \textbf{Output}: `license-plate-int8.onnx`.
        \item \textbf{Result}: Model weights and activations converted from 32-bit floating point to 8-bit integers, effectively compressing the model size by $\approx 75\%$.
    \end{itemize}
\end{frame}

% Verification & Comparison
\begin{frame}{Verification & Comparison}
    \begin{itemize}
        \item \textbf{Inference Setup}: Created a pipeline to run both FP32 and INT8 models using `onnxruntime`.
        \item \textbf{Visualization}: Implemented `draw\_boxes` to overlay detections.
        \item \textbf{Side-by-Side Comparison}:
        \begin{itemize}
            \item Visual validation confirmed that the Quantized INT8 model maintained high accuracy.
            \item Detected bounding boxes and confidence scores were nearly identical to the original FP32 model.
        \end{itemize}
    \end{itemize}
\end{frame}

% The Problem: Execution Provider Error
\begin{frame}{The Challenge: Deployment Failure}
    \begin{itemize}
        \item We exported the model to ONNX using standard dynamic quantization:
        \begin{code}
            yolo export model=yolo11n.pt format=onnx int8=True
        \end{code}
        \item \textbf{Error on Raspberry Pi}:
        \begin{alertblock}{ONNX Runtime Error}
            \texttt{Status: NOT\_IMPLEMENTED : ConvInteger} \\
            \texttt{Could not find an implementation for ConvInteger(10) node with name '/model.0/conv/ConvInteger'}
        \end{alertblock}
        \item \textbf{Reason}: The \texttt{ConvInteger} operator with \textbf{INT8} weights (signed) was not supported by the specific CPU kernel available on the target `onnxruntime` build for aarch64.
    \end{itemize}
\end{frame}

% The Diagnosis
\begin{frame}{Diagnosis & Technical Analysis}
    \begin{itemize}
        \item \textbf{Operator}: \texttt{ConvInteger} performs quantized convolution.
        \item \textbf{Inputs}: Input Data ($x$) and Weights ($w$), plus Zero Points ($x_{zp}, w_{zp}$).
        \item \textbf{Incompatibility}:
        \begin{itemize}
            \item The exported model used \textbf{INT8} (Signed) for weights.
            \item The specific runtime/hardware backend required \textbf{UINT8} (Unsigned) weights for this operator to function correctly.
        \end{itemize}
    \end{itemize}
\end{frame}

% The Fix
\begin{frame}{The Solution: Weight Transformation}
    \begin{itemize}
        \item We developed a custom script \texttt{fix\_quantization.py} to patch the ONNX graph directly.
        \item \textbf{Algorithm}:
        \begin{enumerate}
            \item Load the ONNX model graph.
            \item Iterate through all \texttt{ConvInteger} nodes.
            \item Check if weights are INT8.
            \item \textbf{Shift Weights}: $W_{uint8} = W_{int8} + 128$.
            \item \textbf{Shift Zero Point}: $ZP_{new} = ZP_{old} + 128$.
            \item Update the model initializers in-place.
        \end{enumerate}
    \end{itemize}
\end{frame}

% Code Snippet
\begin{frame}[fragile]{Implementation Details}
    \begin{block}{Python Fix Snippet}
    \begin{verbatim}
# Convert Weights from INT8 to UINT8
w_data_u8 = (w_data.astype(np.int16) + 128).astype(np.uint8)

# Update the Zero Point to match
w_zp_data_u8 = (w_zp_data.astype(np.int16) + 128)
               .astype(np.uint8)

# Replace in ONNX Graph
model.graph.initializer.remove(w_init)
model.graph.initializer.extend([new_w_init])
    \end{verbatim}
    \end{block}
\end{frame}

% Results
\begin{frame}{Results}
    \begin{itemize}
        \item \textbf{Functional Success}: The patched \texttt{fixed\_model.onnx} loads correctly on the Raspberry Pi without errors.
        \item \textbf{Performance}:
        \begin{itemize}
            \item Model Size: Reduced to $\approx 6.5$ MB.
            \item Inference Speed: Real-time capability on CPU.
        \end{itemize}
        \item \textbf{Conclusion}: Direct graph manipulation allowed us to bypass runtime limitations and deploy the quantized model successfully.
    \end{itemize}
\end{frame}

\begin{frame}
    \centering \Huge \textbf{Thank You}
\end{frame}

\end{document}
